\subsection{Avancements}
Le développement d'aujourd'hui a majoritairement de finaliser les services de base et de commencer la télémétrie.
Des outils à utilisation secondaire au projet ont aussi été réalisé avec l'assistance d'intelligence artificielle pour alégé le temps de développement.
Le service responsable des décisions devrait fonctionner au plus tard demain lorsque les simulations GoDot seront terminé.

\subsection{Simulation GoDot (Raphaël Bonhomme)}
La simulation présentait quelques imperfection vis-à-vis l'accélération dans les courbes.
De plus, le robot simulé allait beaucoup trop droit comparativement à celui utilisé dans le projet.
Une solution simple a été d'implémenter des imperfections dans la simulation afin que le modèle subisse des forces extérieurs pour nuire au mouvement.

\subsection{Modification}
Pour optimiser le déploiement des constantes entre la simulation et GoDot, la proposition suivante recommande d'utiliser un service tier qui connectera GoDot et les services embarqués.
L'objectif est de placer les services dans un container \textit{Docker} afin de connecter localement GoDot aux services.